

\section{Conclusion}\label{sec:conclusion} %\marginpar{Please aggiusta le referenze}


In this work, we introduced the notions of convex products (Def.~\ref{def:convprod}) and convex biproduct categories (Def.~\ref{def:convbicat}). We studied their associated monoidal algebra (Prop.~\ref{lemma: copca objects in convbicat}) and established a near analogue of Fox's theorem (Prop.~\ref{prop:A}). We then presented the construction $\stmat{\Cat{C}}$ of stochastic matrices over a $\Cat{PCA}$-enriched category $\Cat{C}$, and showed that $\stmat{\Cat{C}}$ is the free convex biproduct category generated by $\Cat{C}$ (Thm.~\ref{thm:matfree}).

We also developed a syntactic construction: for any category $\Cat{C}$, the category $\CatT{\Cat{C}}$ is obtained by freely adding a monoidal structure equipped with natural monoids and co-pointed convex algebras. We proved that $\CatT{\Cat{C}}$ is the free convex biproduct category over $\Cat{C}$ (Thm.~\ref{thm:syntacticadjunction}). Combining this result with a known theorem (Thm.~\ref{thm:freeenriched}), we derived the isomorphism
$\CatT{\Cat{C}} \cong \stmat{\Cat{C}^+},$
(Cor.~\ref{cor:isotapematrices}), where arrows in $\Cat{C}^+$ are subdistributions over arrows in $\Cat{C}$.

This construction becomes particularly meaningful when $\Cat{C} = \DiagS$, a category of string diagrams. In this case, $\CatT{\DiagS}$ yields the category of probabilistic tape diagrams introduced in~\cite{bonchi2025tapediagramsmonoidalmonads}. The isomorphism
$
\CatT{\DiagS} \cong \stmat{\DiagS^+}
$
thus provides a concrete interpretation of tape diagrams as stochastic matrices whose entries are probability subdistributions over string diagrams.

We applied this framework to probabilistic Boolean circuits with explicit conditioning from~\cite{piedeleu2025boolean}. As illustrated in~\cite[Ex.~30]{bonchi2025tapediagramsmonoidalmonads}, such circuits can be encoded as probabilistic tape diagrams over \emph{partial} Boolean circuits. To obtain a complete axiomatisation of these tapes, we first axiomatised partial Boolean circuits (Thm.~\ref{thm:completenesspartialcircuits}) and then introduced three additional axioms for tapes (Figure~\ref{ax:BooleanTAPES}). Thanks to the isomorphism
$\CatT{\DiagS} \cong \stmat{\DiagS^+}$, the resulting completeness theorem (Thm.~\ref{thm:completenessPBPtapes}) admits a conceptually simple and structurally transparent proof.

\paragraph{Related and future work.} The construction of stochastic matrices over $\Cat{PCA}$-enriched categories is conceptually akin to the classical construction of matrices over categories enriched in commutative monoids~\cite{mac_lane_categories_1978,coecke2017two}. However, while the latter is straightforward enough to be left as an exercise in~\cite{mac_lane_categories_1978}, the former involves significantly more work. In particular, finite biproduct categories enjoy a bijective correspondence between arrows \( f \colon A \oplus B \to C \oplus D \) and quadruples of morphisms \( (f_{A,C}, f_{B,C}, f_{A,D}, f_{B,D}) \), where each \( f_{X,Y} \colon X \to Y \). This correspondence breaks down in convex biproduct categories. 




As illustrated in Section~\ref{sec:probboolcircuits}, our probabilistic tape diagrams are closely related to finitely partially additive categories~\cite{manes}, which are the basis of \emph{effectuses}~\cite{introductioneffectus,chophd}. The main difference is that in an arbitrary finitely partially additive category, one cannot define $\diagp{X}\colon X \to X\oplus X$ unless a notion of scalar is explicitly required. Scalars appear in {effectuses} as arrows of type \( 1 \to 1 \), where \( 1 \) is a distinguished object. As future work, we plan to explore the relationship between our framework and effectuses, which may require considering more structure for tapes, such as their \emph{copy-discard} variants introduced in~\cite{bonchi2025tapediagramsmonoidalmonads}.



%Our probabilistic tape diagrams appear to be closely related to effectuses \cite{introductioneffectus,chophd}. Effectuses are based on finitely \emph{partially additive categories} \cite{manes}, that is, categories enriched over partial commutative monoids, rather than pointed convex algebras.  A precise understanding of the relationship between these two frameworks constitutes an interesting direction for future research.

%This difference is crucial for our approach which is grounded on monoidal algebra. Indeed, in an arbitrary finitely partially additive category, one cannot define $\diagp{X}\colon X \to X\oplus X$ unless some notion of scalar is explicitly required.Scalars appears in \emph{effectuses} \cite{introductioneffectus} as arrows of type $1 \to 1$ where $1$ is a distinguished object. Interestingly in $\CatT{\Cat{C}}$, for some monoidal category $(\Cat{C}, \otimes, 1)$, every  $p\in [0,1]$ defines the arrow $\diagp{1}; (\id{1} \oplus \, \bang{1}) \colon 1 \to 1$. Nevertheless, not all the arrows in $\CatT{\Cat{C}}[1,1]$ are $p\in[0,1]$ (unless $\Cat{C}[1,1]$ is the singleton set): see for instance the syntactic categorty $\Cat{T}(\Diag{B})$ in Section \ref{sec:probboolcircuits}. To make the correspondence with effectuses more precise, one needs to consider more structure for tapes, such as their \emph{copy-discard} variants introduced in \cite{bonchi2025tapediagramsmonoidalmonads}. 


In~\cite{villoria2024enriching}, a general framework is introduced to enrich string diagrams over arbitrary algebraic theories. This approach fundamentally differs from the tape-based construction in~\cite{bonchi2025tapediagramsmonoidalmonads}, in that it does not account for the structure induced by the coproduct \( \oplus \). As a consequence, the correspondence with matrix-based semantics %--central to the tape diagrams and formalised via convex biproducts-- 
does not arise in~\cite{villoria2024enriching}.

We plan to extend probabilistic tape diagrams with traces for \( \oplus \), following the approach developed in~\cite{bonchi2024diagrammatic}. That work showed that a specific property of monoidal traces--called \emph{uniformity}--corresponds to reasoning by invariants in Hoare-style program logics~\cite{hoare1969axiomatic}. Interestingly, as demonstrated in~\cite{jacobs2010coalgebraic}, uniformity holds for traces over \( \oplus \) in \( \KlD \). Preliminary results suggest that, in the probabilistic setting, invariants correspond to sub-martingales, thus offering a promising link between diagrammatic semantics and probabilistic program verification.

Moreover, probabilistic tape diagrams extended with traces are expressive enough to encode probabilistic regular expressions from~\cite{DBLP:conf/lics/RozowskiS24}, which provide a complete axiomatisation for the regular behaviours of generative probabilistic systems--that is, coalgebras in \( \KlD \)~\cite{hasuo2007generic}. We expect that tapes may again offer a more principled and modular axiomatisation, akin to Kleene algebra in~\cite{Kozen94acompleteness}. Interestingly, the completeness proof in~\cite{Kozen94acompleteness} relies heavily on matrix representations over languages; as we hint in Example~\ref{ex:matrices}, our tape formalism naturally supports stochastic matrices over probabilistic languages.

%To conclude, we mention that extending our axiomatisation to probabilistic circuits with $\CBcocopier$ (used in \cite{piedeleu2025boolean} for explicit conditioning) is a challenging future work. Indeed, our proof crucially relies on the well-known isomorphism $\Diag{\mathbb{B}} \cong \Sets_2$, while we are not aware of a similar result for partial functions, necessary to accommodate $\CBcocopier$.

%A further line of investigation concerns the possible connections between probabilistic tape diagrams and the theory of \emph{effectuses} \cite{introductioneffectus}, which offers a general categorical framework for reasoning about probabilistic and quantum systems.

%\emph{Related and future works.} The correspondence between free finite biproduct categories and matrices of arrows of pre-additive categories (see \cite{mac_lane_categories_1978}, \cite{coecke2017two}) relies on the trivial correspondence between functions of the form $f\colon \bigoplus_{i=1}^n U_i \to \bigoplus_{j=1}^m V_j$ and lists of functions $f_{ij}\colon U_i \to V_j$ indexed by $(i,j)\in \{1,\dots,n\} \times \{1,\dots,m\}$. Such correspondence does not hold in convex biproduct categories, which only allow one to build from the $f_{ij}$'s an arrow $f$ that takes into account scaling values $p_{ij}\in [0,1]$ for every $f_{ij}$. Our framework seems promising to be extended to other kinds of enrichment in place of PCA, to obtain diagrams enriched with other algebraic operations such as in \cite{villoria2024enriching}. Another interesting development would be to extend probabilistic tape diagrams with traces for $\oplus$, as done in \cite{bonchi2024diagrammatic} to develop program logic, and to study their algebraic counterpart given by traces of matrices. Such an extension could enable the modelling of stochastic processes. A further line of investigation concerns the possible connections between probabilistic tape diagrams and the theory of \emph{effectuses} \cite{introductioneffectus}, which offers a general categorical framework for reasoning about probabilistic and quantum systems.
